\section{Mean-field approximation: proof of Theorem \ref{2302:thm:mf_approx_lowd_bound}}

\label{2302:sec:app:mf}
\paragraph{Preliminaries} We begin by recalling the results of \cite{mei2019mean}. For any distribution $\mu$ on $\mathbb R^d$, we define the network function
\begin{equation}
    \hat f_{\mu}(\vec{x}) = \int \sigma(\vec{w}^\top \vec{x}) d\mu(\vec{w}).
\end{equation}
It is easy to check that when
\[ \mu = \mu_p(W) := \frac1p \sum_{i = 1}^p \delta_{\vec{w}_i}, \]
then $\hat f_\mu = f_\Theta$, where $f_\Theta$ was defined in \eqref{2302:eq:def:model}. The associated network risk is then given by
\begin{equation}
    \hat{\mathcal{R}}(\mu) = \mathbb E_{\vec{x} \sim \mathcal N(\vec{0}, \sfrac{1}{d} I_d)} \left[\frac12\left( f_{\Theta^\star}(\vec{x}) - \hat f_\mu(\vec{x}) \right)^2 \right].
\end{equation}
Similarly, we can consider the continuous equivalent of the gradient flow equation as follows: 
\begin{equation}
    \varphi(\vec{w}, \mu) = \mathbb E_{\vec{x} \sim \mathcal N(\vec{0}, \sfrac{1}{d} I_d)} \left[\sigma'(\vec{w}^\top \vec{x}) \vec{x} \left( f_{\Theta^\star}(\vec{x}) - \hat f_\mu(\vec{x}) \right) \right].
\end{equation}
Through conservation of matter arguments, the evolution of the empirical measure for particles following the gradient flow equation \eqref{2302:eq:gf} obeys the following partial differential equation:
\begin{equation}
    \label{2302:eq:def:gf_pde}\tag{GF-PDE}
    \partial_t \mu_t = \nabla_{\vec{w}} \cdot \left( \mu_t \varphi( \cdot , \mu_t) \right) 
\end{equation}

\begin{theorem}[\cite{mei2019mean}, Propositions 13-16]\label{2302:thm:mei_mf_bound}
    Let $\mu_0$ be a $K^2/d$-subgaussian distribution, and consider the following two processes:
    \begin{itemize}
        \item the solution $\mu_t$ of \eqref{2302:eq:def:gf_pde} with initial value $\mu_0$,
        \item the solution $W(t)$ of the gradient flow ODE \eqref{2302:eq:gf}, with initial conditions $\vec{w_i}(0) \sim \mu_0$ i.i.d.
    \end{itemize}
    Then for any $T \geq 0$ we have
    \begin{equation}\label{2302:eq:app:mei_mf_bound}
        \sup_{t \in [0, T]} \left|\mathcal{R}(W(t)) - \hat{\mathcal{R}}(\mu_t) \right| \leq \frac{Ce^{CT} \left(\sqrt{\log(pT)} + z \right)}{\sqrt{p}}
    \end{equation}
    with probability at least $1 - e^{-z^2}$.
\end{theorem}

As we have already mentioned in the section devoted to the gradient flow regime, the solution $\Omega(t)$ of \eqref{2302:eq:def:gf_ode} is exactly the overlap matrix of $W(t)$, and hence the term $\mathcal{R}(W(t))$ can be replaced by $\mathcal{R}(\Omega(t))$ in \eqref{2302:eq:app:mei_mf_bound}.

\paragraph{Rotation invariance} The crux of Theorem \ref{2302:thm:mf_approx_lowd_bound} lies in the rotation invariance of $\mu_0$. Indeed, as noticed in \cite{abbe2022merged,hajjar2022symmetries}, such symmetries are conserved throughout the mean-field dynamics.
\begin{proposition}[\cite{hajjar2022symmetries}, Proposition 2.1]\label{2302:prop:app:chizat_symmetry}
Let $U$ be a linear transformation, and assume that the initial measure $\mu_0$, the teacher function $f_\star$ and the data measure $\rho$ are all invariant under $U$. Let $\mu_t$ be the solution to \eqref{2302:eq:def:gf_pde} with initial condition $\mu_0$. Then $\mu_t$ is $U$-invariant for all $t \geq 0$.
\end{proposition}

Write $\mathbb R^d = V^\star \oplus V^\bot$, where $V^\star$ is the span of $W^\star$ and $V^\bot$ is its orthogonal subspace. Proposition \ref{2302:prop:app:chizat_symmetry} then implies immediately that $\mu_t$ is rotation-invariant on $V^\bot$. Hence, if we define the following function:
\begin{equation}
h(\vec{w}) =  \left(\frac{W^\star \vec{w}}d, \lVert \vec{w}^\bot \rVert^2\right)
\end{equation}
where $\vec{w}^\bot$ is the projection of $\vec{w}$ on $V^\bot$, then $\mu_t$ is only determined by its pushforward $\tilde \mu_t = h_\#\mu_t$. 
Conversely, for a set of reduced parameters $(\vec{m}, q)\in\mathbb{R}^d$, and $\vec{g} \in \mathbb{S}^{d-k-1}$, we define
\begin{equation}
    \tilde{\vec{w}} = W^{\star \top}P^{-1}\vec{m} + \sqrt{q} \cdot \vec{g},
\end{equation}
where $\vec{g}$ is a uniform unit vector in $V^\bot$. Then $\mu_t$ is the measure of $\tilde{\vec{w}}$ when $(\vec{m}, q) \sim \tilde \mu_t$ and $\vec{g}\sim \mathrm{Unif}(\mathbb{S}^{d-k-1})$. This allows us to write the reduced equations for $\tilde{\mu_t}$:
\begin{equation}
\label{2302:eq:app:df_pde}\tag{DF-PDE}
    \begin{split}
        \partial_t \tilde\mu_t &= \nabla_{(\vec{m}, q)} \cdot \left( \tilde \mu_t \tilde \varphi( \cdot , \tilde \mu_t) \right) \\
        \tilde \varphi_{\vec{m}}((\vec{m}, q), \tilde \mu) &= \mathbb{E}_{\vec{x}, \vec{g}}\left[ \sigma'(\tilde{\vec{w}}^\top \vec{x}) W^\star \vec{x} \left( f_{\Theta^\star}(\vec{x}) - \tilde f_{\tilde \mu}(\vec{x}) \right) \right] \\
        \tilde \varphi_{q}((\vec{m}, q), \tilde \mu) &= \mathbb{E}_{\vec{x}, \vec{g}}\left[ \sigma'(\tilde{\vec{w}}^\top \vec{x}) \sqrt{q} \vec{g}^\top \vec{x} \left( f_{\Theta^\star}(\vec{x}) -\tilde  f_{\tilde \mu}(\vec{x}) \right) \right]
    \end{split}
\end{equation}
where $\vec{x} = (\vec{z}, \vec{r}) \sim \mathcal{N}(0, \sfrac{1}{d}I_d)$ is a normalised Gaussian vector, $\vec{g} \sim \mathrm{Unif}(\mathbb{S}^{d-k-1})$, and
\begin{equation}
    \tilde f_{\tilde \mu}(\vec{x}) = \int \sigma(\tilde{\vec{w}}^\top \vec{x}) d\tilde\mu(\vec{m}, q) d\nu(\vec{g})
\end{equation}
The associated population risk is now
\begin{equation}
    \tilde{\mathcal{R}}(\tilde\mu) :=  \mathbb E_{\vec{x} \sim \mathcal N(\vec{0}, \sfrac{1}{d} I_d)} \left[\frac12 \left( f_{\Theta^\star}(\vec{x}) - \tilde f_{\tilde\mu}(\vec{x}) \right)^2 \right].
\end{equation}

We have therefore shown the following proposition:
\begin{proposition}
    Assume that $\mu_0$ is rotation-invariant on $V^\bot$. Consider these two processes:
    \begin{itemize}
        \item the solution $\mu_t$ of \eqref{2302:eq:def:gf_pde} with initial value $\mu_0$,
        \item the solution $\tilde\mu_t$ of \eqref{2302:eq:app:df_pde} with initial value $h_{\#}\mu_0$. 
    \end{itemize}
    Then, for all $t \geq 0$, we have
    \begin{equation}
        \hat{\mathcal{R}}(\mu_t) = \tilde{\mathcal{R}}(\tilde\mu_t)
    \end{equation}
\end{proposition}

\paragraph{Back to ODEs} Consider a population of $p$ particles $(\vec{m}, q)$ that evolve according to the following equations:
\begin{equation}
    \label{2302:eq:app:df_ode}\tag{DF-ODE}
    \begin{split}
        \frac{d\vec{m_i}}{dt} &= \tilde \varphi_{\vec{m}}(\vec{m_i}(t), \tilde\mu_{p}(t)) \\
        \frac{dq_i}{dt}& = \tilde \varphi_{q}(q_i(t), \tilde\mu_{p}(t))
    \end{split}
\end{equation}
where $\tilde\mu_p(t)$ is the empirical distribution of the population:
\[ \tilde \mu_p(t) = \frac1p \sum_{i=1}^p \delta_{\vec{m}_i(t), q_i(t)}. \]
Then, by the same arguments as Theorem \ref{2302:thm:mei_mf_bound}, we have:
\begin{proposition}\label{2302:prop:mf_approx_reduced}
    Assume that $\tilde\mu_0$ is a $K^2/d$-subgaussian distribution, and consider the two processes:
    \begin{itemize}
        \item the solution $\tilde\mu_t$ of \eqref{2302:eq:def:gf_pde} with initial value $\tilde\mu_0$,
        \item the solution $\tilde\Theta(t) = (\vec{m_i}(t), q_i(t))_i$ of \eqref{2302:eq:app:df_ode} with initial conditions $\vec{m}_i(0), q_i(0) \sim \tilde \mu_0$ i.i.d.
    \end{itemize}
    Then for all $T \geq 0$, we have
    \begin{equation}
        \sup_{t \in [0, T]} \left|\tilde{\mathcal{R}}(\tilde\Theta(t)) - \tilde{\mathcal{R}}(\tilde\mu_t) \right| \leq \frac{Ce^{CT} \left(\sqrt{\log(pT)} + z \right)}{\sqrt{p}},
    \end{equation}
    with probability at least $1 - e^{-z^2}$, where $\tilde{\mathcal{R}}(\Theta(t)) = \tilde{\mathcal{R}}(\tilde\mu_p(t))$.
\end{proposition}

In conclusion, we have shown the following:
\begin{theorem}\label{2302:thm:app:mf_approx_bound}
    Assume that conditions \ref{2302:assump:activation} and \ref{2302:assump:init} are both satisfied, and define
    \begin{equation}
       \mat{M^0} = \frac{\mat{W}^{0}{\mat{W}^{\star}}^{\top}}{d},\quad  \mat{Q^0} = \frac{\mat{W}^{0}{\mat{W}^{0}}^{\top}}{d},\quad  \vec{q}^0 = \diagarg{Q^0 - M^0 P^{-1} M^{0\top}}.
    \end{equation}
    Consider the two following processes:
    \begin{itemize}
        \item the solution $\Omega(t)$ of \eqref{2302:eq:def:gf_ode} with initial value $\Omega^0$,
        \item the solution $\tilde\Theta(t)$ of \eqref{2302:eq:app:df_ode} with initial value $(M^0, \vec{q}^0)$.
    \end{itemize}
    Then for any $T \geq 0$, we have
    \begin{equation}
        \sup_{t \in [0, T]} \left|\mathcal{R}(\Omega(t)) - \tilde{\mathcal{R}}(\tilde\Theta(t)) \right| \leq \frac{Ce^{CT} \left(\sqrt{\log(pT)} + z \right)}{\sqrt{p}},
    \end{equation}
     with probability at least $1 - e^{-z^2}$.
\end{theorem}

\paragraph{Matching the equations} To show that Theorem \ref{2302:thm:app:mf_approx_bound} implies Theorem \ref{2302:thm:mf_approx_lowd_bound}, we need to show the following:
\begin{itemize}
    \item the update equations for \eqref{2302:eq:def:mf_ode} and \eqref{2302:eq:app:df_ode} are the same
    \item the risk definitions for $\tilde{\mathcal{R}}(\tilde\Theta)$ matches the one from \eqref{2302:eq:def:mf_risk}.
\end{itemize}
We will only show it for $\tilde\varphi_{\vec{m}}$; the rest is done similarly. Expanding the definition, we have
\begin{equation}\label{2302:eq:app:mf_ode_expansion}
    \tilde \varphi_{\vec{m}}((\vec{m}_i, q_i), \tilde \mu_p)_r = \frac1k \sum_{s = 1}^k \langle  \sigma'(\tilde\lambda_i) \lambda^\star_r \sigma(\lambda^\star_s) \rangle - \frac1p \sum_{j=1}^p \langle  \sigma'(\tilde\lambda_i) \lambda^\star_r \sigma(\tilde\lambda_j) \rangle
\end{equation}
where $\tilde \lambda_i = \tilde{\vec{w}}_i^\top \vec{x}$ and $\langle \cdot \rangle$ denotes the expectation with respect to $\vec{x}, \vec{g}$. On the other hand,
\begin{equation}\label{2302:eq:app:ss_ode_expansion}
    \Psi^{(M)}(\Omega) = \frac1k \sum_{s = 1}^k \langle  \sigma'(\lambda_i) \lambda^\star_r \sigma(\lambda^\star_s) \rangle - \frac1p \sum_{j=1}^p \langle  \sigma'(\lambda_i) \lambda^\star_r \sigma(\lambda_j) \rangle,
\end{equation}
without any expectation on $\vec{g}$; hence we only need to match the expressions term by term. The first sums of \eqref{2302:eq:app:mf_ode_expansion} and \eqref{2302:eq:app:ss_ode_expansion} are actually identical (since the marginal distribution of $\lambda_i^\star$ is independent from $\vec{g}$), so we look at the second ones:
\begin{equation}
    \langle  \sigma'(\tilde\lambda_i) \lambda^\star_r \sigma(\tilde\lambda_j) \rangle = \int \mathbb{E}_{x \sim \mathcal{N}(\vec{0}, \sfrac{1}{d}I_d)}\left[ \sigma'(\tilde\lambda_i) \lambda^\star_r \sigma(\tilde\lambda_j) \right] d\nu(\vec{g}_i) d\nu(\vec{g}_j)
\end{equation}
For $i \neq j$ and a given realization of $\vec{g}_i, \vec{g}_j$, the covariance matrix of $(\tilde \lambda_i, \tilde \lambda_j, \lambda^\star_r)$ is
\[ \mathrm{Cov}\left[\tilde \lambda_i, \tilde \lambda_j, \lambda^\star_r\right] = \begin{pmatrix}
    \vec{m}_i^\top P^{-1} \vec{m_i} + q_i & \vec{m}_i^\top P^{-1} \vec{m_j} + \sqrt{q_i q_j} \xi_{ij} & m_{ir} \\
    \vec{m}_i^\top P^{-1} \vec{m_j} + \sqrt{q_i q_j} \xi_{ij} & \vec{m}_i^\top P^{-1} \vec{m_i} + q_i &  m_{jr} \\
    m_{ir} & m_{jr} & p_{rr}
\end{pmatrix} \]
where $\xi_{ij} = \langle \vec{g_i}, \vec{g_j} \rangle$. It is easy to check that this is exactly equal to $\tilde\Omega^{ijr}$, where $\tilde\Omega$ is defined in \eqref{2302:eq:def:overlaps_mf}. Finally, by linearity of expectation, since we never have a three-way correlation term between $\vec{g}_i, \vec{g}_j, \vec{g}_\ell$, we can consider all the $\xi_{ij}$ to be independent. This ends the proof of Theorem \ref{2302:thm:mf_approx_lowd_bound}.
% %%%%%%%%%%%%%%%%%%%%%%%%%%%%%%%%%%%%%%%%%%%%%%%%%%%%%%%%%%%%%%%%%%%%%%%%%%%%%%%
